

\section{Fehlerkette}

\begin{itemize}
	\item Bis 36 Sekunden nach dem Start zeigte die Rakete normales Flugverhalten.
	\item Danach Fehler des Back-Up Trägheitsnavigationssystems und ganz kurz darauf Fehler des aktiven Trägheitsnavigationssystems.
	\item Das aktive Trägheitsnavigationssystem sendet einen Fehlercode an den On Board Computer und schaltet sich ab.
	\item Versuch auf das Back-Up umzuschalten, aber dieses hat sich ja bereits einen Zyklus (72 ms) vorher verabschiedet.
	\item Der On Board Rechner interpretiert die erhaltenen Fehlercodes als Flugbahndaten und führt daher komplett falsche Berechnungen durch.
	\item Steuerdüsen der Rakete werden bis zum Maximum ausgeschwenkt, um die vermeintlich falsche Flugposition zu korrigieren.
	\item Dies führt zu enormen Belastungen der Verbindungen zwischen dem Booster und der Hauptstufe der Rakete. Diese Verbindungen gehen dadurch zu Bruch.
	\item Die Rakete wird gesprengt, wie es für dieses Ereignis vorgesehen ist, um zu verhindern, dass die Rakete als ganzes auf die Erde einschlägt.e
\end{itemize}
